\clearpage

% 2. OBJETIVOS
\section{Objetivos}

Objetivos claros fornecem direção e propósito ao Clube de Matemática. Ao defini-los desde
o início, o grupo estabelece um entendimento comum sobre o que pretende alcançar e como o
sucesso será medido.

\subsection{Objetivo Principal}

\begin{infobox}
  \textbf{Desenvolver uma base sólida de raciocínio matemático e lógico entre os cadetes
  da 42 São Paulo, contribuindo diretamente para a permanência, o pertencimento e o
  desempenho no curso.}
\end{infobox}

Para alcançar esse objetivo, o clube atuará sobre três pilares:

\begin{enumerate}
  \item Fortalecimento de fundamentos matemáticos essenciais para a programação;
  \item Criação de um espaço de comunidade e apoio mútuo entre cadetes;
  \item Redução da taxa de abandono causada por dificuldades lógico-matemáticas.
\end{enumerate}

\subsection{Objetivos Específicos}

\subsubsection{Desenvolver Competências Matemáticas Essenciais}

O clube focará no desenvolvimento progressivo de conhecimentos matemáticos práticos e
aplicáveis ao contexto da 42. Por meio de sessões estruturadas, exercícios colaborativos e
desafios progressivos, os membros construirão uma base sólida nas áreas definidas no
Capítulo~5 deste regimento.

O sucesso será medido por:

\begin{itemize}
  \item Realização de ao menos uma sessão por semana de forma contínua;
  \item Participação ativa de ao menos 8 cadetes por sessão;
  \item Relatos positivos de membros sobre impacto na confiança e na permanência.
\end{itemize}

\subsubsection{Reduzir Barreiras de Entrada no Cursus}

Muitos cadetes abandonam a 42 nas primeiras semanas por dificuldades com lógica, raciocínio
abstrato ou matemática básica. O clube oferecerá suporte específico nesses pontos críticos,
criando uma rede de apoio para quem está travado e se sente sozinho.

\subsubsection{Criar Senso de Comunidade}

Além do conteúdo matemático, o clube funcionará como um ponto de encontro regular --- um
grupo menor na comunidade maior da 42, onde o cadete tem rostos conhecidos, pessoas
que torcem por ele e um motivo concreto para aparecer no campus.

\subsection{Objetivos de Consolidação}

\subsubsection{Produzir Proposta Pedagógica Formal}

Após os primeiros ciclos de atividade, o clube documentará sua metodologia, resultados e
impacto observado em um documento formal a ser apresentado à equipe pedagógica da 42 São
Paulo. O objetivo é institucionalizar o grupo, demonstrando de forma concreta sua
contribuição para a retenção e o desempenho dos cadetes no cursus.

\subsection{Realismo e Adaptabilidade}

Os objetivos deste regimento são ambiciosos, mas alcançáveis para um grupo motivado de
cadetes. Com sessões regulares, participação ativa e liderança compartilhada, é realista
construir uma comunidade matemática relevante dentro da 42 em 6 a 12 meses. Se necessário,
o grupo revisará periodicamente suas metas e condições de funcionamento.
